% ===========================================================================
% MODELO DE RELATORIO  ·  ENE0046 Laboratorio de Eletronica  ·  2026/2
%
% Como usar
%   1. Preencha o bloco IDENTIFICACAO logo abaixo.
%   2. Deixe comentados os autores que nao existirem. Um autor fica centralizado,
%      dois dividem a linha, tres ocupam tercos.
%   3. Substitua o texto de cada secao pelo seu. As orientacoes que estao la
%      dizem o que se espera em cada uma.
%   4. Compile com pdflatex duas vezes, ou com latexmk -pdf.
%
% Estrutura
%   Roteiro 1 nao tem bancada: use so a subsecao 4.1 e a 5.1, e compare o calculo
%   analitico com a simulacao de projeto.
%   Roteiros 2 a 7 usam a estrutura inteira.
%
% O arquivo estilo-relatorio.tex e a imagem logo-unb.png devem estar na mesma
% pasta deste arquivo.
% ===========================================================================
\documentclass[11pt,a4paper]{article}
% !TeX root = modelo-relatorio.tex
% Formatacao do relatorio. Compile o modelo-relatorio.tex, nao este arquivo.
\usepackage[T1]{fontenc}
\usepackage[utf8]{inputenc}
\usepackage{babel}
\usepackage[a4paper,top=1.6cm,bottom=1.8cm,left=2.1cm,right=2.1cm]{geometry}
\usepackage[default]{lato}
\usepackage{amsmath,amssymb}
\usepackage{siunitx}
\usepackage{graphicx}
\usepackage{booktabs,array,longtable,colortbl}
\usepackage{enumitem}
\usepackage{tcolorbox}
\usepackage{xcolor}
\usepackage{tikz}
\usepackage{titlesec}
\usepackage{fancyhdr}
\usepackage{caption}
\usepackage{etoolbox}
\usepackage{placeins}
\usepackage{float}
\usepackage[colorlinks=true,linkcolor=azul,urlcolor=azul,citecolor=azul]{hyperref}
\tcbuselibrary{skins,breakable}

\definecolor{azul}{HTML}{0A4B8C}
\definecolor{azulclaro}{HTML}{E8F0F8}
\definecolor{verde}{HTML}{0B7A4B}
\definecolor{cinza}{HTML}{5B6672}
\definecolor{cinzaclaro}{HTML}{F4F6F8}
\definecolor{vinho}{HTML}{B04A4A}
\definecolor{vinhoclaro}{HTML}{FBF2F2}

\sisetup{output-decimal-marker={,}, per-mode=symbol}
\setlength{\parindent}{0pt}
\setlength{\parskip}{6pt plus 1pt}
\setlist{itemsep=3pt, topsep=4pt, parsep=0pt, leftmargin=1.3em}
\renewcommand{\arraystretch}{1.25}
\linespread{1.06}

% ---- secoes numeradas com marcador colorido -------------------------------
\titleformat{\section}[hang]
  {\Large\bfseries\color{azul}}
  {\raisebox{-1pt}{\tikz\node[fill=azul,text=white,rounded corners=2pt,
     inner sep=3pt,minimum width=20pt,align=center]{\bfseries\thesection};}}
  {10pt}{}
\titlespacing*{\section}{0pt}{16pt plus 3pt}{6pt}
\titleformat{\subsection}{\large\bfseries\color{cinza}}{\thesubsection}{8pt}{}
\titlespacing*{\subsection}{0pt}{11pt plus 2pt}{3pt}

% ---- legendas -------------------------------------------------------------
\captionsetup{font=small,labelfont={bf,color=azul},labelsep=period,skip=6pt}

% ---- identificacao --------------------------------------------------------
\newcommand{\autorA}{}\newcommand{\matA}{}
\newcommand{\autorB}{}\newcommand{\matB}{}
\newcommand{\autorC}{}\newcommand{\matC}{}
\newcommand{\autorum}[2]{\renewcommand{\autorA}{#1}\renewcommand{\matA}{#2}}
\newcommand{\autordois}[2]{\renewcommand{\autorB}{#1}\renewcommand{\matB}{#2}}
\newcommand{\autortres}[2]{\renewcommand{\autorC}{#1}\renewcommand{\matC}{#2}}

\newcommand{\disciplina}{CIC0236 · Teleinformática e Redes 2}
\newcommand{\semestre}{2026/2}
\newcommand{\turma}{}
\newcommand{\roteironum}{}
\newcommand{\roteirotema}{}
\newcommand{\aturma}[1]{\renewcommand{\turma}{#1}}
\newcommand{\roteiro}[2]{\renewcommand{\roteironum}{#1}\renewcommand{\roteirotema}{#2}}

\newcommand{\umautor}[2]{%
  \begin{minipage}[t]{\dimexpr\linewidth-4pt\relax}\centering
    \parbox[t][30pt][t]{\linewidth}{\centering
      {\fontsize{10.5}{12.5}\selectfont #1}}\\[-2pt]
    {\color{cinza}\fontsize{9}{11}\selectfont\texttt{#2}}
  \end{minipage}}

\newcounter{nautores}
\newcommand{\blocoautores}{%
  \setcounter{nautores}{0}%
  \ifdefempty{\autorA}{}{\stepcounter{nautores}}%
  \ifdefempty{\autorB}{}{\stepcounter{nautores}}%
  \ifdefempty{\autorC}{}{\stepcounter{nautores}}%
  \ifnum\value{nautores}=0
    \begin{center}\color{vinho}\bfseries
      Preencha ao menos um autor no preâmbulo do arquivo.
    \end{center}
  \else
    \begin{center}
    \ifcase\value{nautores}\or
      % --- um autor: ocupa o centro
      \begin{minipage}{0.5\textwidth}\umautor{\autorA}{\matA}\end{minipage}%
    \or
      % --- dois autores: metades
      \begin{minipage}{0.42\textwidth}\umautor{\autorA}{\matA}\end{minipage}\hspace{0.04\textwidth}%
      \begin{minipage}{0.42\textwidth}\umautor{\autorB}{\matB}\end{minipage}%
    \or
      % --- tres autores: tercos
      \begin{minipage}{0.31\textwidth}\umautor{\autorA}{\matA}\end{minipage}\hspace{0.015\textwidth}%
      \begin{minipage}{0.31\textwidth}\umautor{\autorB}{\matB}\end{minipage}\hspace{0.015\textwidth}%
      \begin{minipage}{0.31\textwidth}\umautor{\autorC}{\matC}\end{minipage}%
    \fi
    \end{center}
  \fi}

% ---- cabecalho do relatorio -----------------------------------------------
\newcommand{\capa}{%
\begin{tcolorbox}[colback=azul,colframe=azul,boxrule=0pt,arc=4pt,
  left=16pt,right=16pt,top=13pt,bottom=13pt]
\color{white}
\begin{minipage}{0.76\textwidth}
  {\fontsize{9}{11}\selectfont\textbf{\MakeUppercase{\disciplina{} · \semestre{}}}}\\[6pt]
  {\fontsize{22}{25}\selectfont\bfseries Relatório da Entrega \roteironum}\\[3pt]
  {\fontsize{12.5}{15}\selectfont \roteirotema}
\end{minipage}\hfill
\begin{minipage}{0.2\textwidth}\raggedleft
  \colorbox{white}{\hspace{2pt}\raisebox{-2pt}{\includegraphics[height=30pt]{logo-unb.png}}\hspace{2pt}}
\end{minipage}
\end{tcolorbox}
\vspace{10pt}
\blocoautores
\ifdefempty{\turma}{}{%
  \begin{center}{\color{cinza}\fontsize{9}{11}\selectfont Turma \turma}\end{center}}
\vspace{4pt}
{\color{azulclaro}\hrule height 2pt}
\vspace{2pt}}

% ---- caixas auxiliares ----------------------------------------------------
\newtcolorbox{formula}{colback=azulclaro,colframe=azul,boxrule=0.8pt,arc=3pt,
  left=12pt,right=12pt,top=9pt,bottom=9pt}
\newtcolorbox{nota}[1]{colback=cinzaclaro,colframe=cinza,boxrule=0.8pt,arc=3pt,
  left=10pt,right=10pt,top=7pt,bottom=7pt,
  title={\textbf{#1}},fonttitle=\small\bfseries,coltitle=white,colbacktitle=cinza}
\newenvironment{resumo}
  {\vspace{6pt}\begin{center}\begin{minipage}{0.92\textwidth}
   \begin{tcolorbox}[colback=cinzaclaro,colframe=cinzaclaro,boxrule=0pt,arc=3pt,
     left=14pt,right=14pt,top=10pt,bottom=10pt,
     borderline west={2.5pt}{0pt}{azul}]
   {\color{azul}\fontsize{8}{10}\selectfont\textbf{RESUMO}}\\[4pt]
   \small}
  {\end{tcolorbox}\end{minipage}\end{center}\vspace{2pt}}

\newtcolorbox{espec}{colback=azulclaro,colframe=azul,boxrule=0.8pt,arc=3pt,
  left=12pt,right=12pt,top=9pt,bottom=9pt,
  title={\textbf{Especificação da matrícula}},fonttitle=\bfseries,
  coltitle=white,colbacktitle=azul}

% ---- rodape ---------------------------------------------------------------
\pagestyle{fancy}\fancyhf{}
\renewcommand{\headrulewidth}{0pt}
\renewcommand{\footrulewidth}{0.4pt}
\fancyfoot[L]{\color{cinza}\fontsize{8}{10}\selectfont \disciplina{} · \semestre{} · Entrega \roteironum}
\fancyfoot[R]{\color{cinza}\fontsize{8}{10}\selectfont \thepage}

\newlist{referencias}{itemize}{1}
\setlist[referencias]{label={}, leftmargin=1.1em, itemindent=-1.1em,
                     itemsep=3pt, topsep=2pt, parsep=0pt}

\newcommand{\mat}[1]{\texttt{\textbf{#1}}}


% --------------------------- IDENTIFICACAO ---------------------------------
\roteiro{2}{Janela Deslizante: GBN e SR}
\aturma{01}

\autorum{Enzo Morrone de Azevedo Muniz}{221035442}
\autordois{Luca Valderramos Cirino}{211066140}
\autortres{Raul Myron Silva Amorim}{200049712}
% ---------------------------------------------------------------------------

\begin{document}
\capa

% Entrega e avaliação
%
% Relatório objetivo em PDF (máximo 8 páginas), contendo obrigatoriamente:

\begin{resumo}
Este relatório apresenta a implementação e a avaliação de protocolos de transferência confiável sobre UDP
utilizando janelas deslizantes. Foram comparadas as estratégias Go-Back-N (GBN) e Selective Repeat (SR),
com janelas de 4, 8 e 16 segmentos, tendo o Stop-and-Wait (SAW) como referência. Os testes foram realizados
na transferência de um arquivo de \SI{1}{\mebi\byte}, dividido em segmentos de 512 bytes, em cenários sem
perdas e com perdas artificiais de 5\% e 10\%. A integridade dos arquivos recebidos foi confirmada pelo checksum
MD5 nos três protocolos. Sem perdas, o aumento da janela elevou o desempenho da transferência. Sob perdas, o SR
apresentou maior eficiência e menor tempo de transferência, pois retransmite somente os segmentos perdidos, enquanto
o GBN retransmite também segmentos posteriores à falha.
\end{resumo}

\section{Introdução}

Com a expansão da internet para além dos PCs, se tornou necessário desenvolver aplicações com tempo de resposta mínimo.
Isso deve-se ao celular, hoje o principal meio de se conectar na internet, que exige protocolos rápidos - devido as exigências dos usuários -
e confiáveis - devido a perda de pacotes esperada das comunicações sem fio.

Dessa forma, um pode pensar que o protocolo TCP é ideal, porém, devido a perda de pacotes, o mecanismo de confirmar pacotes com ACK
pode cause retransmissões excessivas, diminuindo a largura de banda da comunicação. Já o UDP também não é ideal, visto que não garante
a comunicação confiável, apesar de permitir uma largura de banda melhor.

Assim, este relatório é parte 2 de 3, visando criar um novo protocolo que mistura os pontos fortes do UDP e TCP afim de obter
o melhor dos dois mundos - confiável e rápido. Neste relatório 1, usando o UDP como base, é implementado uma melhoria sobre o
\textit{Stop-and-Wait} (SAW) do último relatório, o uso de janelas de transmissão nas técnicas \textit{Go-Back-N} (GBN) e \textit{Selective Repeat} (SR).

Em ambos os protocolos, é utilizado uma janela deslizante, onde é mantido a lista dos \textit{n} últimos REQs enviados sem ACK.
Assim, ao contrário do SAW, ocorre mais de uma transferência de cada vez. A diferença está em como cada um age em caso de um \textit{timeout}.

No \textit{Go-Back-N}, se qualquer segmento sofrer \textit{timeout}, retransmite ele e todos os seguintes; já no \textit{Selective Repeat},
retransmite apenas o segmento com \textit{timeout} (sem retransmitir os demais).

Assim, espera-se que ambas as novas técnicas obtenham um desempenho melhor que o SAW devido ao uso de janelas de transmissão. Além disso,
em caso de situações com uma grande perda de pacotes, espera-se que o SR funcione melhor de o GBN, ao evitar requisitar de novo pacotes que
já foram obtidos com sucesso. Para isso, deve ser implementado um mecanismo de perda de pacotes artificial para testes.

Este relatório está dividido em 3 secções principais: na implementação é dito como foi feito o código que implementa o protocolo;
nos resultados, tem-se a verificação da confiabilidade com o checksum e dados de throughput e retransmissão da transmissão;
e por último na análise, os dados dos resultados são comentados.

\section{Implementação}

%%%%%%%%%%%%%%%%%%%%%%%%%%%%%%%%%%%%%%%%%%%%%%%%%%%%%%%%%%%%%%%%%%%%%%%%%%%%%%%%
% Descrição da implementação — para cada protocolo (GBN e SR), apresente um fluxograma ou máquina de estados descrevendo o comportamento da janela, o critério de retransmissão e o tratamento de segmentos fora de ordem (SR)
%%%%%%%%%%%%%%%%%%%%%%%%%%%%%%%%%%%%%%%%%%%%%%%%%%%%%%%%%%%%%%%%%%%%%%%%%%%%%%%%

%%%%%%%%%%%%%%%%%%%%%%%%%%%%%%%%%%%%%%%%%%%%%%%%%%%%%%%%%%%%%%%%%%%%%%%%%%%%%%%%
% Método de simulação de perda — descreva qual abordagem o grupo escolheu (descartar DATA ou não enviar ACK), justifique a escolha e explique como o campo simulated_loss foi registrado no CSV
%%%%%%%%%%%%%%%%%%%%%%%%%%%%%%%%%%%%%%%%%%%%%%%%%%%%%%%%%%%%%%%%%%%%%%%%%%%%%%%%


\section{Resultados}

\subsection{Verificação}
\begin{table}[H]
\begin{center}
\begin{tabular}{@{}lr@{}}
\toprule
	\textbf{Protocolo} & \textbf{Checksum MD5 calculado} \\
\midrule
	SAW & 564f9d97cd203fee214569a3de1b54aa \\
	GBN & 564f9d97cd203fee214569a3de1b54aa \\
	SR  & 564f9d97cd203fee214569a3de1b54aa \\
\bottomrule
\end{tabular}
\end{center}
\caption{Dados obtidos dos protocolos SAW, GBN e SR para a transmissão do arquivo grupo12 small}
	\label{tab:throughput}
\end{table}

\subsection{Dados}
Cada entrada da tabela está no formado Throughput médio|Retransmissões totais |Tempo total

Durante as medições, o RTT médio variou entre \SI{30}{\milli\second} a \SI{45}{\milli\second}.

Nas simulações de perda, todas as medições usaram a mesma seed para a geração de número aleatório. Apesar disso, há diferença no número de retransmissões
devido a perda de pacotes não artificiais e as diferenças dos protocolos no tratamento de timeouts

\begin{table}[H]
\begin{center}
\begin{tabular}{@{}lrrr@{}}
\toprule
	\textbf{Protocolo} & \textbf{Sem Perda} & \textbf{5\% de perda} & \textbf{10\% de perda} \\
\midrule

	SAW    & \SI{130.92}{\kilo\bit\per\second}  & \SI{105.53}{\kilo\bit\per\second} & \SI{104.15}{\kilo\bit\per\second} \\
	GBN-4  & \SI{395.63}{\kilo\bit\per\second}  & \SI{107.46}{\kilo\bit\per\second} & \SI{66.12}{\kilo\bit\per\second} \\
	GBN-8  & \SI{759.18}{\kilo\bit\per\second}  & \SI{145.63}{\kilo\bit\per\second} & \SI{88.63}{\kilo\bit\per\second} \\
	GBN-16 & \SI{1497.06}{\kilo\bit\per\second} & \SI{194.51}{\kilo\bit\per\second} & \SI{121.41}{\kilo\bit\per\second} \\
	SR-4   & \SI{375.31}{\kilo\bit\per\second}  & \SI{243.00}{\kilo\bit\per\second} & \SI{167.35}{\kilo\bit\per\second} \\
	SR-8   & \SI{889.34}{\kilo\bit\per\second}  & \SI{488.41}{\kilo\bit\per\second} & \SI{329.40}{\kilo\bit\per\second} \\
	SR-16  & \SI{1648.73}{\kilo\bit\per\second} & \SI{933.47}{\kilo\bit\per\second} & \SI{642.78}{\kilo\bit\per\second} \\

\bottomrule
\end{tabular}
\end{center}
	\caption{Throughput médio obtido dos protocolos SAW, GBN e SR para a transmissão do arquivo grupo12 small de \SI{1}{\mebi\byte} de tamanho}
	\label{tab:throughput}
\end{table}

\begin{table}[H]
\begin{center}
\begin{tabular}{@{}lrrr@{}}
\toprule
	\textbf{Protocolo} & \textbf{Sem Perda} & \textbf{5\% de perda} & \textbf{10\% de perda} \\
\midrule

	SAW    & 0 & 122 & 250 \\
	GBN-4  & 0 & 294 & 566 \\
	GBN-8  & 0 & 458 & 777 \\
	GBN-16 & 0 & 789 & 1408 \\
	SR-4   & 0 & 122 & 249 \\
	SR-8   & 0 & 122 & 249 \\
	SR-16  & 0 & 122 & 249 \\

\bottomrule
\end{tabular}
\end{center}
	\caption{Retransmissões totais obtidas dos protocolos SAW, GBN e SR para a transmissão do arquivo grupo12 small de \SI{1}{\mebi\byte} de tamanho}
	\label{tab:throughput}
\end{table}

\begin{table}[H]
\begin{center}
\begin{tabular}{@{}lrrr@{}}
\toprule
	\textbf{Protocolo} & \textbf{Sem Perda} & \textbf{5\% de perda} & \textbf{10\% de perda} \\
\midrule

	SAW    & \SI{64.07}{\second} & \SI{79.49}{\second} & \SI{80.54}{\second} \\
	GBN-4  & \SI{21.20}{\second} & \SI{78.06}{\second} & \SI{126.88}{\second} \\
	GBN-8  & \SI{11.05}{\second} & \SI{57.60}{\second} & \SI{94.64}{\second} \\
	GBN-16 & \SI{5.60}{\second}  & \SI{43.13}{\second} & \SI{69.09}{\second} \\
	SR-4   & \SI{22.35}{\second} & \SI{34.52}{\second} & \SI{50.13}{\second} \\
	SR-8   & \SI{9.43}{\second}  & \SI{17.18}{\second} & \SI{25.47}{\second} \\
	SR-16  & \SI{5.09}{\second}  & \SI{8.99}{\second}  & \SI{13.05}{\second} \\

\bottomrule
\end{tabular}
\end{center}
	\caption{Tempo total obtido dos protocolos SAW, GBN e SR para a transmissão do arquivo grupo12 small de \SI{1}{\mebi\byte} de tamanho}
	\label{tab:throughput}
\end{table}

%%%%%%%%%%%%%%%%%%%%%%%%%%%%%%%%%%%%%%%%%%%%%%%%%%%%%%%%%%%%%%%%%%%%%%%%%%%%%%%%
% Gráficos por segmento — throughput e número de retransmissões ao longo da transferência, para pelo menos um cenário de perda, comparando GBN e SR com a mesma janela
%%%%%%%%%%%%%%%%%%%%%%%%%%%%%%%%%%%%%%%%%%%%%%%%%%%%%%%%%%%%%%%%%%%%%%%%%%%%%%%%

\begin{figure}[H]
\centering
	% \includegraphics[width=0.8\linewidth]{}
\caption{Throughput e número de retransmissões ao longo da transferência}
\label{fig:graf_rtt}
\end{figure}


\section{Analise}

%%%%%%%%%%%%%%%%%%%%%%%%%%%%%%%%%%%%%%%%%%%%%%%%%%%%%%%%%%%%%%%%%%%%%%%%%%%%%%%%
% Análise obrigatória: a) Qual o impacto do tamanho da janela no throughput de GBN e SR sem perda? A partir de qual tamanho de janela o throughput se estabiliza? Por quê? 
%%%%%%%%%%%%%%%%%%%%%%%%%%%%%%%%%%%%%%%%%%%%%%%%%%%%%%%%%%%%%%%%%%%%%%%%%%%%%%%%
Ao analisar os resultados, no caso sem perda, é possível observar uma relação diretamente proporcional entre o tamanho da janela e o throughput médio.

%%%%%%%%%%%%%%%%%%%%%%%%%%%%%%%%%%%%%%%%%%%%%%%%%%%%%%%%%%%%%%%%%%%%%%%%%%%%%%%%
% b) Com 10% de perda e janela 16, quantos segmentos o GBN retransmite para cada erro? E o SR? Mostre isso com dados do CSV. 
%%%%%%%%%%%%%%%%%%%%%%%%%%%%%%%%%%%%%%%%%%%%%%%%%%%%%%%%%%%%%%%%%%%%%%%%%%%%%%%%

%%%%%%%%%%%%%%%%%%%%%%%%%%%%%%%%%%%%%%%%%%%%%%%%%%%%%%%%%%%%%%%%%%%%%%%%%%%%%%%%
% c) Por que o GBN sofre mais que o SR com alta taxa de perda? Responda com os números da tabela comparativa.
%%%%%%%%%%%%%%%%%%%%%%%%%%%%%%%%%%%%%%%%%%%%%%%%%%%%%%%%%%%%%%%%%%%%%%%%%%%%%%%%


\section{Conclusão}

Neste trabalho foi implementada uma extensão do protocolo confiável desenvolvido na etapa anterior, substituindo o mecanismo Stop-and-Wait por protocolos baseados em janelas deslizantes. Foram consideradas as estratégias Go-Back-N e Selective Repeat, que permitem manter múltiplas requisições de segmentos em andamento e tratam as perdas por meio de timeouts e retransmissões.

A implementação utiliza o UDP como mecanismo de transporte, acrescentando controle de sequência, confirmação dos segmentos recebidos e reconstrução ordenada do arquivo. Também foi incluída uma forma de simular perdas durante a transferência, permitindo avaliar o comportamento dos protocolos em diferentes condições de comunicação. A verificação por checksum MD5 permite confirmar que, apesar das perdas e retransmissões, o arquivo recebido permanece íntegro.

De modo geral, o uso de janelas tende a aproveitar melhor a capacidade da conexão do que o Stop-and-Wait, enquanto a diferença entre Go-Back-N e Selective Repeat se torna especialmente relevante quando ocorrem perdas. A comparação quantitativa entre as janelas e os cenários de perda será complementada na análise dos resultados apresentados neste relatório.


\section{Referências}

\begin{referencias}
\item James F. Kurose \& Keith W. Ross, \emph{Redes de Computadores e a Internet - Uma nova Abordagem}, 8a Edição, Pearson Education / Makron Books.

\item Andrew S. Tanenbaum, \emph{Redes de Computadores}, Pearson Universidades, 5a Edição, 2011
\end{referencias}

\end{document}
